% ============================================================
%  SECCIÓN — CÁLCULOS ESTRUCTURALES (MURO EN MÉNSULA)
%  Armado del fuste y de la zapata — ELU Código Estructural 2021
% ============================================================
\section{Cálculos estructurales del muro en ménsula}
\label{sec:muro_estructural}

El dimensionamiento estructural tiene como objetivo verificar que la sección
transversal del muro soporta las acciones de cálculo bajo Estado Límite Último
(ELU) de acuerdo con el Código Estructural (CE, RD~470/2021). Se estudian tres
secciones críticas: la \textbf{base del alzado} (máximo momento flector y
cortante), el \textbf{talón de la zapata} (ménsula cargada por las tierras) y
la \textbf{puntera de la zapata} (ménsula ascendente por la presión de
contacto).

\begin{figure}[H]
  \centering
  \includegraphics[width=0.80\textwidth]{figures/muro/imagen_armaduras}
  \caption{Esquema de armaduras del muro en ménsula.}
  \label{fig:imagen_armaduras}
\end{figure}

% -----------------------------------------------------------
\subsection{Normativa y coeficientes de cálculo}
% -----------------------------------------------------------

El dimensionamiento estructural se realiza según el
\textbf{Código Estructural} (CE, Real Decreto 470/2021). Se adopta
coeficiente parcial único $\gamma_f = 1{,}50$ para todas las acciones
(tierras + agua + sobrecarga), criterio conservador habitual en muros
de contención.

\begin{table}[H]
  \centering
  \caption{Propiedades de cálculo de los materiales}
  \label{tab:mat_calculo}
  \begin{tabular}{@{}llcc@{}}
    \toprule
    \textbf{Material} & \textbf{Parámetro} & \textbf{Característico}
    & \textbf{Cálculo} \\
    \midrule
    HA-30 & $f_{ck}$                         & 30~MPa  & — \\
          & $f_{cd} = f_{ck}/\gamma_c$        & —       & 20~MPa \\
          & $f_{ctm,fl} = 0{,}30\,f_{ck}^{2/3}$ & — & 2{,}90~MPa \\
    B~500~SD & $f_{yk}$                       & 500~MPa & — \\
             & $f_{yd} = f_{yk}/\gamma_s$     & —       & 434~MPa \\
    \midrule
    \multicolumn{4}{l}{$\gamma_c = 1{,}50$;\quad $\gamma_s = 1{,}15$;\quad
      Recubrimiento nominal: $r_{\mathrm{nom}} = 50$~mm} \\
    \bottomrule
  \end{tabular}
\end{table}

% -----------------------------------------------------------
\subsection{Acciones sobre el alzado}
\label{subsec:acciones_fuste}
% -----------------------------------------------------------

El alzado actúa como \textbf{ménsula vertical} empotrada en la zapata
(altura $H = 7{,}30$~m). Las acciones horizontales se calculan de forma
análoga al análisis global de estabilidad
(sección~\ref{subsec:empujes_rankine_opc}), pero adaptadas a la
altura del alzado (sin la zapata), se puede comprobar en el anexo de cálculos manuales.
El empuje del agua actúa desde el N.F.\ ($-2{,}50$~m) hasta la base del
alzado ($h_w = 7{,}30 - 2{,}50 = 4{,}80$~m).

\begin{table}[H]
  \centering
  \caption{Componentes del empuje horizontal sobre el alzado (por ml)}
  \label{tab:empuje_fuste}
  \begin{tabular}{@{}lcccc@{}}
    \toprule
    \textbf{Componente} & \textbf{Forma}
    & $E$ \textbf{(Tn/ml)} & $x$ \textbf{(m)} & $M$ \textbf{(Tn·m/ml)} \\
    \midrule
    $E_1$ — Arena no sat.\ ($h_1=2{,}50$~m) & Triáng.\ & 1{,}97 & 4{,}96 &  9{,}77 \\
    $E_2$ — Arena sat.\  ($h_2=1{,}00$~m)  & Mixta     & 1{,}75 & 4{,}28 &  7{,}49 \\
    $E_{q,1}$ — Sobrecarga en arena         & Rect.\    & 1{,}21 & 5{,}55 &  6{,}72 \\
    $E_{q,2}$ — Sobrecarga en zona saturada & Rect.\    & 1{,}93 & 1{,}90 &  3{,}67 \\
    $E_w$ — Agua ($h_w=4{,}80$~m)           & Triáng.\  & 11{,}52 & 2{,}40 & 27{,}65 \\
    $E_3$ — Arcilla ($e_2$ decreciente)     & Triáng.\  &  2{,}26 & 1{,}26 &  2{,}86 \\
    \midrule
    \textbf{TOTAL} & & $\mathbf{E_T = 20{,}64}$ & — & $\mathbf{M_T = 58{,}16}$ \\
    \bottomrule
  \end{tabular}
\end{table}

Altura de aplicación de la resultante sobre la base del alzado:
\[
  h_{ef} = \frac{M_T}{E_T} = \frac{58{,}16}{20{,}64} = 2{,}82~\text{m}
\]

% -----------------------------------------------------------
\subsection{Combinación de carga ELU — alzado}
\label{subsec:elu_fuste}
% -----------------------------------------------------------

Aplicando el coeficiente $\gamma_f = 1{,}50$ a todas las acciones:

\begin{align}
  M_d &= \gamma_f \cdot M_T = 1{,}50 \times 58{,}16
       = \resultado{87{,}24~\text{Tn·m/ml}}
  \label{eq:Md_fuste} \\[4pt]
  V_d &= \gamma_f \cdot E_T = 1{,}50 \times 20{,}64
       = \resultado{30{,}96~\text{Tn/ml} = 303{,}72~\text{kN/m}}
  \label{eq:Vd_fuste}
\end{align}

% -----------------------------------------------------------
\subsection{Armado del alzado}
\label{subsec:armado_fuste}
% -----------------------------------------------------------

\subsubsection{Altura útil en la sección de empotramiento}

La sección más solicitada es la \textbf{base del alzado}
($e_2 = 0{,}75$~m). Con $r_{\mathrm{nom}} = 50$~mm:
\[
  d = e_2 - r_{\mathrm{nom}} - \frac{\phi_{\mathrm{main}}}{2}
    = 750 - 50 = \resultado{d = 0{,}70~\text{m}}
\]

\subsubsection{Armadura principal — trasdós del alzado}

\paragraph{Comprobación de momento límite (sección rectangular):}
\begin{align}
  M_{\mathrm{lim}} &= 0{,}372\,f_{cd}\,b\,d^2
    = 0{,}372 \times 20\,000 \times 1 \times 0{,}70^2 \notag \\
    &= 3\,645~\text{kN·m/m} = 371{,}6~\text{Tn·m/ml}
    \gg M_d = 87{,}24~\text{Tn·m/ml} \quad\checkmark
\end{align}

No se necesita armadura de compresión. Armadura simple.

\paragraph{Área de acero necesaria:}
\begin{equation}
  A_{s,\mathrm{req}} = \frac{M_d}{0{,}85\,d\,f_{yd}}
    = \frac{87{,}24}{0{,}85 \times 0{,}70 \times 434 \times 10^2}
    = \resultado{3{,}37~\text{cm}^2/\text{m}}
  \label{eq:As_fuste_calc}
\end{equation}

\paragraph{Cuantía mínima (CE 2021, Art.\ 42.3.5):}

Se aplica la cuantía mínima de flexión para muros:
\begin{equation}
  A_{s,\min} = \frac{A_c \cdot f_{ctm,fl}}{4{,}8 \cdot f_{yd}}
    = \frac{1 \times 0{,}75 \times 2{,}90}{4{,}8 \times 434}
    = \resultado{10{,}44~\text{cm}^2/\text{m}}
  \label{eq:As_min_fuste}
\end{equation}

Como $A_{s,\mathrm{req}} = 3{,}37~\text{cm}^2/\text{m} < A_{s,\min} = 10{,}44~\text{cm}^2/\text{m}$,
\textbf{rigen las cuantías mínimas del CE}.

\paragraph{Elección de barras — alzado:}

Se selecciona $\phi 16$ ($A_{\phi 16} = 2{,}01~\text{cm}^2$):
\[
  n = \frac{A_{s,\min}}{A_{\phi 16}} = \frac{10{,}44}{2{,}01} = 5{,}19
  \;\longrightarrow\; 6~\text{barras/m}
  \qquad
  s = \frac{100}{6} = 16{,}7~\text{cm}
\]

Comprobación de separación según CE: $s = 16{,}7~\text{cm}$;
$s_{\min} = \max(20~\text{mm};\,\phi;\,\mathrm{TMA}+5) = \max(20;16;25) = 25$~mm
$= 2{,}5~\text{cm} \ll 16{,}7~\text{cm}$\checkmark;
$s \leq 35~\text{cm}$\checkmark.

\begin{resultadobox}[Armadura principal alzado — trasdós]
  \[
    6\,\phi 16 \;\text{c}/16{,}7~\text{cm}
    \quad\Rightarrow\quad
    A_s = 6 \times 2{,}01 = \resultado{12{,}06~\text{cm}^2/\text{m}}
    > A_{s,\min} = 10{,}44~\text{cm}^2/\text{m} \quad\checkmark
  \]
\end{resultadobox}

% -----------------------------------------------------------
\subsection{Comprobación a cortante — alzado}
\label{subsec:cortante_fuste}
% -----------------------------------------------------------

El cortante de cálculo en la base del alzado es:
\[
  V_{Ed} = 30{,}96~\text{Tn/ml} \times 9{,}81~\text{kN/Tn}
         = \resultado{303{,}72~\text{kN/m}}
\]

\paragraph{Resistencia a cortante sin armadura ($V_{Rd,c}$):}
\begin{align*}
  C_{Rd,c} &= \frac{0{,}18}{\gamma_c} = \frac{0{,}18}{1{,}5} = 0{,}12 \\[4pt]
  k &= 1 + \sqrt{\frac{200}{d[\mathrm{mm}]}}
     = 1 + \sqrt{\frac{200}{700}} = 1{,}535 \leq 2 \quad\checkmark \\[4pt]
  \rho_l &= \frac{A_{sl}}{b_w \cdot d}
          = \frac{12{,}06 \times 10^{-4}~\text{m}^2/\text{m}}{0{,}75 \times 0{,}70}
          = 0{,}0023 \leq 0{,}02 \quad\checkmark \\[4pt]
  V_{Rd,c} &= \left[C_{Rd,c}\cdot k\cdot(100\,\rho_l\,f_{ck})^{1/3}\right]b_w\,d \\
           &= \left[0{,}12\times 1{,}535\times(100\times 0{,}0023\times 30)^{1/3}\right]
             \times 750 \times 700 \\
           &= \left[0{,}12\times 1{,}535\times 1{,}905\right]\times 525\,000
            = 0{,}351\times 525\,000 = \resultado{184\,275~\text{N/m} \approx 184{,}5~\text{kN/m}}
\end{align*}

Valor mínimo:
\[
  V_{Rd,c,\min} = 0{,}035\,k^{3/2}\,f_{ck}^{1/2}\,b_w\,d
    = 0{,}035\times 1{,}535^{3/2}\times\sqrt{30}\times 750\times 700
    = \resultado{191{,}6~\text{kN/m}}
\]

\begin{resultadobox}[Comprobación a cortante — alzado]
  \[
    V_{Ed} = 303{,}72~\text{kN/m} \;>\; V_{Rd,c} = 184{,}5~\text{kN/m}
  \]
  \begin{center}
    \textbf{SE NECESITA ARMADURA DE CORTANTE (estribos)}
  \end{center}
\end{resultadobox}

\subsubsection{Dimensionamiento de estribos}

\paragraph{Comprobación de las bielas de compresión:}
\begin{equation}
  V_{Rd,\max} = 0{,}27\!\left(1-\frac{f_{ck}}{250}\right)f_{cd}\,b_w\,d
    = 0{,}27\times 0{,}88\times 20\times 750\times 700
    = 2\,494\,800~\text{N/m} \gg V_{Ed} \quad\checkmark
\end{equation}

\paragraph{Cuantía de armadura transversal necesaria:}

La resistencia del acero de cortante se limita a $f_{ywd} = 400$~MPa según CE:
\begin{equation}
  \frac{A_{sw}}{s} = \frac{V_{Ed}}{0{,}9\,d\,f_{ywd}}
    = \frac{303\,720~\text{N/m}}{0{,}9\times 700\times 400}
    = \resultado{1\,205~\text{mm}^2/\text{m}}
  \label{eq:Asw_s}
\end{equation}

\paragraph{Cuantía mínima de armadura transversal (CE):}
\[
  A_{sw,\min} = \frac{0{,}08\sqrt{f_{ck}}}{f_{yk}}\,b_w \cdot s
  = \frac{0{,}08\sqrt{30}}{500}\times 750 \times 1\,000
  = 657~\text{mm}^2/\text{m}
  < \frac{A_{sw}}{s} = 1\,205~\text{mm}^2/\text{m} \quad\checkmark
\]

\paragraph{Elección de estribos — 3 ramas $\phi 10$:}

Área por estribo: $A_{\mathrm{cerco}} = 3\times\frac{\pi\times 10^2}{4}
= 235{,}6~\text{mm}^2$.

Separación longitudinal:
\[
  s_l = \frac{A_{\mathrm{cerco}}}{A_{sw}/s}
      = \frac{235{,}6}{1\,205} = 0{,}1955~\text{m}
  \;\longrightarrow\; \resultado{s_l = 195~\text{mm}}
\]
Comprobación: $s_l = 195~\text{mm} \leq s_{l,\max} = 0{,}75\,d = 525~\text{mm}$\checkmark.

Separación transversal entre ramas:
\[
  s_t = \frac{750 - 2\times 50}{2} = 325~\text{mm}
  \leq s_{t,\max} = 0{,}75\,d = 525~\text{mm} \quad\checkmark
\]

\begin{resultadobox}[Armadura de cortante alzado]
  \[
    3~\text{ramas}\;\phi 10 \;\text{c}/195~\text{mm}
  \]
  \[
    \frac{A_{sw}}{s} = \frac{235{,}6}{195}\times 1\,000
    = \resultado{1\,208~\text{mm}^2/\text{m}} > 1\,205~\text{mm}^2/\text{m}\;\checkmark
  \]
\end{resultadobox}

% -----------------------------------------------------------
\subsection{Cálculo de la zapata}
\label{subsec:calculo_zapata}
% -----------------------------------------------------------

\subsubsection{Distribución de tensiones de contacto}

Para el armado de la zapata, la tensión de contacto se obtiene
a partir de la resultante de fuerzas verticales y su excentricidad.
\textbf{No se descuentan los momentos volcadores}: se usa directamente
$N_T = 68{,}80$~Tn/ml y $M_s = 203{,}96$~Tn·m/ml (momentos
estabilizadores respecto a la puntera).

Posición de la resultante $N_T$ desde el borde de la puntera:
\begin{equation}
  e = \frac{M_s}{N_T} = \frac{203{,}96}{68{,}80} = 2{,}964~\text{m}
\end{equation}

Excentricidad respecto al centro de la zapata:
\begin{equation}
  e' = \frac{B}{2} - e = \frac{5{,}60}{2} - 2{,}964 = -0{,}164~\text{m}
  \quad\text{(hacia el talón)}
\end{equation}

Distribución de presiones (trapezoidal):
\begin{align}
  \sigma &= \frac{N_T}{B}\!\left(1 \pm \frac{6\,e'}{B}\right)
           = \frac{68{,}80}{5{,}60}\!\left(1 \pm \frac{6\times 0{,}164}{5{,}60}\right)
           = 12{,}29\,(1 \pm 0{,}176) \notag \\[4pt]
  \sigma_{\max} &= 12{,}29 \times 1{,}176
    = \resultado{14{,}45~\text{Tn/m}^2} \quad\text{(bajo el talón)} \\
  \sigma_{\min} &= 12{,}29 \times 0{,}824
    = \resultado{10{,}13~\text{Tn/m}^2} \quad\text{(bajo la puntera)}
\end{align}

Ambas positivas: no hay despegue del cimiento\checkmark.

\subsubsection{Armadura del talón (cara superior de la zapata)}
\label{subsubsec:armado_talon}

El talón ($B_{\mathrm{talón}} = 4{,}05$~m) trabaja como ménsula
empotrada en la cara del alzado. La presión en la sección de empotramiento
(interpolando a distancia $4{,}05$~m del extremo del talón):

\begin{equation}
  \sigma_{\mathrm{encast}} = \sigma_{\max} - (\sigma_{\max}-\sigma_{\min})
    \cdot\frac{B_{\mathrm{talón}}}{B}
  = 14{,}45 - 4{,}32\times\frac{4{,}05}{5{,}60}
  = \resultado{11{,}33~\text{Tn/m}^2}
\end{equation}

Momento flector en la sección de empotramiento del talón:
\begin{align}
  M_{\mathrm{talón}} &= \frac{\sigma_{\mathrm{encast}}\cdot L_t^2}{2}
    + \frac{(\sigma_{\max} - \sigma_{\mathrm{encast}})\cdot L_t^2}{3} \notag \\
  &= \frac{11{,}33\times 4{,}05^2}{2}
    + \frac{(14{,}45-11{,}33)\times 4{,}05^2}{3}
  = 92{,}91 + 17{,}06 = \resultado{109{,}97~\text{Tn·m/ml}}
\end{align}

Momento de cálculo ELU:
\begin{equation}
  M_{d,\mathrm{talón}} = \gamma_f \cdot M_{\mathrm{talón}}
    = 1{,}50 \times 109{,}97 = \resultado{164{,}96~\text{Tn·m/ml}}
  \label{eq:Md_talon}
\end{equation}

Comprobación de momento límite (canto útil $d_z = h_z - r = 0{,}80 - 0{,}05 = 0{,}75$~m):
\[
  M_{\mathrm{lim}} = 0{,}372\,f_{cd}\,b\,d_z^2
    = 0{,}372\times 20\,000\times 1\times 0{,}75^2
    = 4\,185~\text{kN·m/m} = 426{,}6~\text{Tn·m/ml} \gg 164{,}96 \quad\checkmark
\]

Área de acero necesaria:
\begin{equation}
  A_{s,\mathrm{req}}^{\mathrm{talón}}
    = \frac{M_{d,\mathrm{talón}}}{0{,}85\,d_z\,f_{yd}}
    = \frac{164{,}96}{0{,}85 \times 0{,}75 \times 434 \times 10^2}
    = \resultado{5{,}96~\text{cm}^2/\text{m}}
  \label{eq:As_talon_calc}
\end{equation}

Cuantía mínima (CE, Art.\ 42.3.5, con $h_z = 0{,}80$~m):
\begin{equation}
  A_{s,\min}^{\mathrm{zapata}}
    = \frac{A_c\,f_{ctm,fl}}{4{,}8\,f_{yd}}
    = \frac{1\times 0{,}80\times 2{,}90}{4{,}8\times 434}
    = \resultado{11{,}13~\text{cm}^2/\text{m}}
  \label{eq:As_min_zapata}
\end{equation}

Como $A_{s,\mathrm{req}} = 5{,}96 < A_{s,\min} = 11{,}13$~cm²/m,
\textbf{rigen las cuantías mínimas del CE}.

\begin{resultadobox}[Armadura superior talón]
  \[
    6\,\phi 16 \;\text{c}/16{,}7~\text{cm}
    \quad\Rightarrow\quad
    A_s = 6 \times 2{,}01 = \resultado{12{,}06~\text{cm}^2/\text{m}}
    > 11{,}13~\text{cm}^2/\text{m} \quad\checkmark
  \]
\end{resultadobox}

\subsubsection{Armadura de la puntera (cara inferior de la zapata)}
\label{subsubsec:armado_puntera}

La puntera ($B_{\mathrm{punta}} = 0{,}80$~m) trabaja como ménsula ascendente.
Presión en la sección de empotramiento (a $0{,}80$~m del extremo de puntera):
\begin{equation}
  \sigma_{\mathrm{encast,punta}} = \sigma_{\min}
    + (\sigma_{\max} - \sigma_{\min})\cdot\frac{0{,}80}{B}
  = 10{,}13 + 4{,}32\times\frac{0{,}80}{5{,}60}
  = \resultado{10{,}75~\text{Tn/m}^2}
\end{equation}

Momento flector en la sección de empotramiento de la puntera:
\begin{align}
  M_{\mathrm{punta}} &= \frac{\sigma_{\mathrm{encast,punta}}\cdot L_p^2}{2}
    + \frac{(\sigma_{\mathrm{encast,punta}} - \sigma_{\min})\cdot L_p^2}{3} \notag \\
  &= \frac{10{,}75\times 0{,}80^2}{2}
    + \frac{(10{,}75-10{,}13)\times 0{,}80^2}{3}
  = 3{,}44 + 0{,}13 = \resultado{3{,}57~\text{Tn·m/ml}}
\end{align}

Momento de cálculo ELU:
\begin{equation}
  M_{d,\mathrm{punta}} = 1{,}50 \times 3{,}57 = \resultado{5{,}36~\text{Tn·m/ml}}
  \label{eq:Md_punta}
\end{equation}

Área de acero necesaria:
\[
  A_{s,\mathrm{req}}^{\mathrm{punta}}
    = \frac{5{,}36}{0{,}85\times 0{,}75\times 434\times 10^2}
    = \resultado{0{,}19~\text{cm}^2/\text{m}}
    \ll A_{s,\min} = 11{,}13~\text{cm}^2/\text{m}
\]

Rige la cuantía mínima. Se adopta la misma solución que en el talón:

\begin{resultadobox}[Armadura inferior puntera]
  \[
    6\,\phi 16 \;\text{c}/16{,}7~\text{cm}
    \quad\Rightarrow\quad
    A_s = \resultado{12{,}06~\text{cm}^2/\text{m}} > 11{,}13~\text{cm}^2/\text{m}\;\checkmark
  \]
\end{resultadobox}

\subsubsection{Armadura en espera (arranque del alzado)}
\label{subsubsec:armado_espera}

La armadura en espera debe ser capaz de transmitir el mismo esfuerzo
que la armadura del alzado:
\[
  A_{s,\mathrm{espera}} = A_{s,\mathrm{fuste}}
  = 6\,\phi 16~\text{c}/16{,}7~\text{cm}
  \;\Rightarrow\; 12{,}06~\text{cm}^2/\text{m}
\]

Longitud de solape (Tabla~2.8.1.2.b CE, $f_{ck}=30~\text{MPa}$,
$f_{yk}=500~\text{MPa}$, supuesto~1):
\[
  l_b = 580~\text{mm}
\]

Como el canto de la zapata ($h_z = 800$~mm) es menor que la longitud de
solape neta ($l_{b,\text{neta}} = 812$~mm), se propone \textbf{doblar la
armadura en espera en forma de patilla en la base de la zapata}:
\[
  l_b(\text{patilla}) = 580 \times 0{,}70 \times 1{,}00 = 406~\text{mm}
  \;\longrightarrow\; l_s = 1{,}4\times 406 \approx 570~\text{mm} \approx 57~\text{cm}
\]
\[
  l_{\mathrm{patilla}} = 10\,\phi = 10\times 16 = 160~\text{mm} = 16~\text{cm}
\]

% -----------------------------------------------------------
\subsection{Resumen de armaduras — cuadro general}
\label{subsec:resumen_armaduras}
% -----------------------------------------------------------

\begin{table}[H]
  \centering
  \caption{Resumen de armaduras del muro en ménsula (por ml de muro)}
  \label{tab:armaduras_muro}
  \resizebox{\textwidth}{!}{%
  \begin{tabular}{@{}llccccc@{}}
    \toprule
    \textbf{Elemento} & \textbf{Posición}
    & \textbf{Barra} & \textbf{Esp. (cm)}
    & $A_s$ \textbf{(cm²/m)}
    & \textbf{Long. barra (m)} & \textbf{kg/ml} \\
    \midrule
    \multirow{3}{*}{Alzado}
      & Trasdós (principal)      & $\phi 16$ & 16{,}7 & 12{,}06 & 7{,}50 &  71{,}5 \\
      & Cortante (estribos)      & $\phi 10$ & 19{,}5 & — (12{,}08\,mm²/mm) & — & 30{,}4 \\
      & Dist.\ horiz.\ (cara)   & $\phi 10$ & 15{,}0 &  5{,}24 & 1{,}00 & 32{,}5 \\
    \midrule
    \multirow{3}{*}{Zapata}
      & Sup.\ talón              & $\phi 16$ & 16{,}7 & 12{,}06 & 4{,}85 &  46{,}1 \\
      & Inf.\ puntera            & $\phi 16$ & 16{,}7 & 12{,}06 & 1{,}60 &  15{,}2 \\
      & Transversal (distribuc.) & $\phi 12$ & 20{,}0 &  5{,}65 & 1{,}00 &  11{,}3 \\
    \midrule
    \multicolumn{5}{l}{\textbf{TOTAL por ml de muro}} & & \textbf{$\approx$207~kg} \\
    \multicolumn{5}{l}{\textbf{TOTAL muro 150~ml}} & &
      \textbf{$\approx$31\,050~kg $\approx$ 31~Tm} \\
    \bottomrule
  \end{tabular}}
\end{table}

\begin{table}[H]
  \centering
  \caption{Verificación final — Estado Límite Último (CE 2021)}
  \label{tab:verificacion_elu}
  \resizebox{\textwidth}{!}{%
  \begin{tabular}{@{}lcccc@{}}
    \toprule
    \textbf{Sección} & $M_d$ \textbf{(Tn·m/ml)} & $A_{s,\mathrm{req}}$
    \textbf{(cm²/m)} & $A_{s,\mathrm{prov}}$ \textbf{(cm²/m)} & \textbf{Estado} \\
    \midrule
    Alzado — base (trasdós)  & 87{,}2 &  3{,}37 (mín: 10{,}44) & 12{,}06 & \checkmark\,(mín.) \\
    Zapata — talón (sup.)    & 165{,}0 & 5{,}96 (mín: 11{,}13) & 12{,}06 & \checkmark\,(mín.) \\
    Zapata — puntera (inf.)  & 5{,}4  & 0{,}19 (mín: 11{,}13) & 12{,}06 & \checkmark\,(mín.) \\
    \midrule
    \multicolumn{2}{l}{Cortante alzado} & $V_{Ed}=303{,}7$~kN/m &
    $V_{Rd,c}=184{,}5$~kN/m & 3$\phi$10\,c/195\,mm \\
    \bottomrule
  \end{tabular}}
\end{table}
